\chapter*{Introduction}
\addcontentsline{toc}{chapter}{Introduction}

Ce rapport présente la conception et l'intégration d'une chaîne de perception
neuromorphique pour la navigation autonome d'un robot mobile LIMO équipé d'une
caméra DAVIS346.

Le cœur du projet est l'implémentation d'une méthode de détection d'objets en
mouvement basée sur la compensation d'ego-mouvement par IMU dans un flux
d'événements. Cette approche est ensuite raccordée à Nav2 afin d'injecter des
obstacles dynamiques dans la prise de décision de navigation.

Le document est organisé comme suit:
\begin{itemize}
\item le chapitre 1 introduit le matériel et les fondements théoriques;
\item le chapitre 2 détaille la méthodologie d'implémentation de l'article
retenu et son adaptation ROS 2;
\item le chapitre 4 présente le protocole de test et l'analyse des résultats;
\item la conclusion synthétise les acquis et les perspectives.
\end{itemize}
